\documentclass{exam}

\newif\ifA
\newif\ifkey
\newif\ifprintselected   % required by exam_config's \ansbox

%%%%%%%%%%%%%%%%%%%%%%%%%%%%%%%%%%%%%%%%%%%%%%%%%%%%%%
%%% SET THESE IF VARIABLES
%%%%%%%%%%%%%%%%%%%%%%%%%%%%%%%%%%%%%%%%%%%%%%%%%%%%%%
% Lab name -- labs are named by the lesson they belong to, never by a
% sequential lab number (Neil, 2026-10-12): this is the L19 Lab.
\def\labLesson{19}

% is this the key or not?
\ifdefined\iskey
	\ifnum\iskey=1
		\keytrue
	\else
		\keyfalse
	\fi
\else
	%% Pick one manually
	\keyfalse % if this is NOT the key
%	\keytrue     % if this IS the key
\fi

% set up the header and footer - change for every term
\ifdefined\thisTermIn
	\def\thisterm{\thisTermIn}
\else
	\def\thisterm{Fall 2026}
\fi
%
\def\thisclass{ECE 444}

%%%%%%%%%%%%%%%%%%%%%%%%%%%%%%%%%%%%%%%%%%%%%%%%%%%%%%
%%% END SET THESE IF VARIABLES
%%%%%%%%%%%%%%%%%%%%%%%%%%%%%%%%%%%%%%%%%%%%%%%%%%%%%%
% packages and shortcuts
\input{myPackages_gr}
\usepackage[hidelinks, linkcolor=red]{hyperref}
\input{myShortcuts}
\usepackage{listings}
\lstset{basicstyle=\ttfamily\small, columns=fullflexible, keepspaces=true,
	showstringspaces=false, breaklines=true}
\setlength{\parindent}{0pt}
\setlength{\parskip}{6pt}

%% exam class customization
\input{exam_config}
% ECE 444 document fonts.
%
% Body text is Barlow (SIL OFL; the faces and the license are in
% latex/fonts/Barlow/). Barlow is not in TeX Live, so it is loaded as an
% OpenType face through fontspec -- which means the document must be built
% with lualatex, not pdflatex. build_practice.sh and build_project.sh do; the
% lab packets build in ECE444-faculty from its vendored copy of this file.
%
% Compiling with pdflatex anyway is not an error: the guard below falls back
% to the LaTeX default so the source still compiles under any engine and can
% be syntax-checked. The resulting PDF just is not the one we ship.
%
% Paths are relative to latex/, which is where the build scripts run
% (each does cd "$(dirname "$0")"). Loading by path rather than by installed
% family name means no system font install and no fontconfig, so a fresh
% clone builds a byte-comparable PDF anywhere.
%
% MATH IS DELIBERATELY LEFT IN COMPUTER MODERN. Barlow has no math companion,
% and CM math under sans body text is already the house pairing: the decks set
% words in Source Sans Pro and let MathJax render symbols in its CM-derived
% face. Moving math to a sans as well would mean unicode-math plus a sans math
% font (Fira Math being the plausible one) and re-checking every equation in
% the course for symbol coverage. That is a separate decision, not a side
% effect of this one.

\usepackage{iftex}
\ifLuaTeX
	\usepackage{fontspec}
	\defaultfontfeatures{Ligatures=TeX}
	\setmainfont{Barlow}[
		Path           = fonts/Barlow/,
		Extension      = .ttf,
		UprightFont    = *-Regular,
		ItalicFont     = *-Italic,
		BoldFont       = *-Bold,
		BoldItalicFont = *-BoldItalic,
	]
	\setsansfont{Barlow}[
		Path           = fonts/Barlow/,
		Extension      = .ttf,
		UprightFont    = *-Regular,
		ItalicFont     = *-Italic,
		BoldFont       = *-SemiBold,
		BoldItalicFont = *-BoldItalic,
	]
	% microtype earns its place here specifically because of Barlow. Its
	% glyphs are narrower than Computer Modern's, so justified prose sets
	% tighter and TeX has to squeeze interword space to fit; two paragraphs
	% in the L02/L03 sets went overfull on the first Barlow build. Protrusion
	% and font expansion absorb that -- both are lualatex features, so this
	% belongs inside the guard with everything else.
	\usepackage{microtype}

	% Mono stays Latin Modern: Barlow ships no monospaced face, and verbatim
	% and code need one that actually is fixed-pitch.
\fi
   % Barlow body text (needs lualatex)
\printselectedfalse      % full worked solutions in the key, not answers-only
\CorrectChoiceEmphasis{\color{red}}
\SolutionEmphasis{\color{red}}
\renewcommand\questionlabel{\textbf{\thequestion.}}
\renewcommand{\thepartno}{\alph{partno}}
\renewcommand{\partlabel}{(\thepartno)}

% Fillable table cell: blank in the student copy, the expected value in red on
% the KEY. Fixed-width centered columns keep the blank cells writable.
\newcolumntype{K}[1]{>{\centering\arraybackslash}p{#1}}
\newcommand{\kv}[1]{\ifkey{\color{red}#1}\fi\rule[-0.10in]{0pt}{0.34in}}

% print the answers if this is the key
\ifkey
	\printanswers % if this IS the key
\else
	\noprintanswers  %if this is NOT the key
\fi

% print the header and footer
% need to print key on top if this is the key
\ifkey
	\firstpageheader{}{\color{red}{ KEY}}{}
	\runningheader{\thisclass\hspace{1pt} -- \thisterm }{L\labLesson\ Lab\hspace{1pt}\color{red}{ KEY}}{\thepage\hspace{1pt} of\hspace{1pt} \numpages}
	\firstpagefooter{}{}{}
	\runningfooter{}{}{}
\else
	\firstpageheader{}{}{}
	\runningheader{\thisclass \hspace{1pt} -- \thisterm }{L\labLesson\ Lab\hspace{1pt}}{\thepage\hspace{1pt} of\hspace{1pt} \numpages}
	\firstpagefooter{}{}{}
	\runningfooter{}{}{}
\fi

\runningheadrule


\begin{document}
\pagestyle{headandfoot}   % exam class


%%%%%%%%%%%%%%%%%%%%%%%%%%%%%%%%%%%%%%%%%%%%%%%%%%%%%%
% title block
%%%%%%%%%%%%%%%%%%%%%%%%%%%%%%%%%%%%%%%%%%%%%%%%%%%%%%

\begin{center}
USAF ACADEMY \\
DEPARTMENT OF ELECTRICAL AND COMPUTER ENGINEERING \\
LESSON \labLesson\ LAB - Fox Hunt \\
\textbf{Lab Packet} \\
\thisterm \\[4mm]
\end{center}

\vspace{-4pt}
\textbf{Names:} \rule[-2pt]{2.2in}{0.4pt}\ \ \rule[-2pt]{2.2in}{0.4pt} \hfill
\textbf{Section:} \rule[-2pt]{0.7in}{0.4pt}

\vspace{4pt}

\fullwidth{\textbf{Learning Objective 3.5: } I can implement beam steering on the
ADALM-PHASER and verify the steered pattern against theory.}

\vspace{2pt}

One HB100 is on, inside one of five numbered boxes downrange of a taped firing
line. Your team moves its kit along the line, stopping at station marks every
$0.5\ \m$ with the kit squared to the tape and broadside pointed straight
downrange, and finds the live box from the bearings and levels it reads. This
packet is the L\labLesson\ Lab, in three parts, and your team turns in one copy,
complete, before Lesson 20. \textbf{Part 1}, before the lab period, is the error
budget, a station plan, and your fox-hunt script, tested on a supplied set of
readings; you need all three at the bench. \textbf{Part 2} is your team's
measurements, recorded in class. \textbf{Part 3}, after the lab, compares your
fix with the box the instructor reveals.

\textbf{Grading.} One submission per team, counted toward every partner's
engagement grade and graded as a genuine, documented attempt.

The array is the PHASER's: $N = 8$ elements, E1 to E8, on a $d = 14\ \text{mm}$
pitch. Use $c = 3.00\times 10^{8}\ \m/\text{s}$ and the HB100's nominal
$10.525\ \ghz$. The range frame puts $x$ along the tape, with station 0 at its
middle, and $y$ straight downrange. Every Angle $\theta$ in this packet
has your reference-shot sign applied, so a positive angle points toward the
$+x$ end of the tape, and a bearing from station $x_i$ is the line
\eq{ x = x_i + y\tanp{\theta_i}. }

\begin{questions}

\fullwidth{\large\textbf{Part 1: Pre-Lab} \normalsize(before the lab period)}

%%%%%%%%%%%%%%%%%%%%%%%%%%%%%%%%%%%%%%%%%%%%%%%%%%%
\question \textbf{The bearing error budget.} The \textbf{Angle} readout, set to Peak, is the steer angle of the
largest sample in the beam sweep, and the sweep steps the progressive phase one
ADAR1000 LSB, $360\mydeg/2^7 = 2.8125\mydeg$, at a time.

\begin{parts}

	\part Find $kd$ in degrees, the sweep's step in $\sin\theta$, and the step in
	angle at each Angle in the table. Then give the largest error the grid
	alone can cause at each angle.

	\begin{center}
	\renewcommand{\arraystretch}{1.25}
	\begin{tabular}{| l | K{1.1in} | K{1.1in} | K{1.1in} |}
	\hline
	Angle & $0\mydeg$ & $30\mydeg$ & $45\mydeg$ \\\hline
	Grid step & \kv{$0.91\mydeg$} & \kv{$1.05\mydeg$} & \kv{$1.29\mydeg$} \\\hline
	Largest grid error & \kv{$0.46\mydeg$} & \kv{$0.53\mydeg$} & \kv{$0.64\mydeg$} \\\hline
	\end{tabular}
	\end{center}

		\begin{solution}[1.2in]
		\eq{ \lambda &= \frac{3.00\times 10^{8}}{10.525\times 10^{9}} = 28.50\ \text{mm}, \qquad
		     kd = 360\mydeg\ \frac{14}{28.50} = 176.8\mydeg \\
		     \delta(\sin\theta) &= \frac{2.8125\mydeg}{176.8\mydeg} = 0.0159, \qquad
		     \delta\theta = \frac{0.0159}{\cosp{\theta}}\ \text{rad} }
		The step widens as $1/\cos\theta$, and Angle can miss the source by
		half a step.
		\end{solution}

	\part On the range, every kit is squared to the tape with a straightedge,
	which holds its heading to about $1.0\mydeg$. Take $1.0\mydeg$ for
	multipath and $0.2\mydeg$ for noise. Combine these with the grid term at
	$30\mydeg$ to find the bearing uncertainty $\sigma_\beta$. Then repeat for a
	kit aimed at a point by eye, which holds the heading only to about
	$1.5\mydeg$.

	\begin{center}
	$\sigma_\beta$ (squared to tape) = \ansbox[1.1in]{$1.5\mydeg$}
	\hspace{0.4in}
	$\sigma_\beta$ (aimed by eye) = \ansbox[1.1in]{$1.9\mydeg$}
	\end{center}

		\begin{solution}[1.0in]
		The terms are independent, so they add as a root sum of squares:
		\eq{ \sigma_\beta &= \sqrt{0.53^2 + 1.0^2 + 1.0^2 + 0.2^2} = 1.5\mydeg \\
		     &= \sqrt{0.53^2 + 1.5^2 + 1.0^2 + 0.2^2} = 1.9\mydeg }
		Squaring to one tape also makes much of the heading error common to every
		station, which rotates all the lines together instead of scattering them.
		\end{solution}

	\part A teammate sets the kit on the floor, $0.75\ \m$ below the boxes, at a
	station $2\ \m$ from a box that is $30\mydeg$ off broadside in azimuth. What
	does Angle read, and which way does the error always pull?

	\begin{solution}[1.2in]
	The box is $\tan^{-1}(0.75/2) = 20.6\mydeg$ above the array's plane. A line
	array measures the cone angle:
	\eq{ \sin\theta_\text{meas} = \sinp{30\mydeg}\cosp{20.6\mydeg} = 0.468, \qquad \theta_\text{meas} = 27.9\mydeg }
	The error always pulls the reading toward broadside, since $\cos(\text{el}) \le 1$.
	Keep the kit on a stand or table at the boxes' height.
	\end{solution}

\end{parts}

%%%%%%%%%%%%%%%%%%%%%%%%%%%%%%%%%%%%%%%%%%%%%%%%%%%
\newpage
\question \textbf{The range.} The firing line is $3.0\ \m$ of tape with stations
at $x = -1.5, -1.0, \ldots, +1.5\ \m$. The box field is $2.0\ \m$ wide,
centered on station 0, and runs from $1.0$ to $2.0\ \m$ downrange; the boxes are
at least $0.75\ \m$ apart. Take $\sigma_\beta = 1.5\mydeg$ in this question.

\begin{parts}

	\part Find the across-range error of one bearing line at each depth.

	\begin{center}
	\renewcommand{\arraystretch}{1.25}
	\begin{tabular}{| l | K{1.1in} | K{1.1in} | K{1.1in} |}
	\hline
	Depth $y$ & $1.0\ \m$ & $1.5\ \m$ & $2.0\ \m$ \\\hline
	Across-range error & \kv{$0.03\ \m$} & \kv{$0.04\ \m$} & \kv{$0.05\ \m$} \\\hline
	\end{tabular}
	\end{center}

		\begin{solution}[0.5in]
		$\delta_x = y\tanp{\sigma_\beta}$; at $2.0\ \m$, $2.0\tanp{1.5\mydeg} = 0.05\ \m$.
		\end{solution}

	\part Every bearing on a firing line looks downrange, so depth is the weak
	direction. Two stations a baseline $B$ apart fix a box at depth $y$ in depth
	to about $y^2\sigma_\beta/B$, with $\sigma_\beta$ in radians. Find the depth
	error for a box $2.0\ \m$ out at each baseline.

	\begin{center}
	\renewcommand{\arraystretch}{1.25}
	\begin{tabular}{| l | K{1.1in} | K{1.1in} | K{1.1in} |}
	\hline
	Baseline $B$ & $0.5\ \m$ & $1.0\ \m$ & $2.0\ \m$ \\\hline
	Depth error & \kv{$0.21\ \m$} & \kv{$0.10\ \m$} & \kv{$0.05\ \m$} \\\hline
	\end{tabular}
	\end{center}

		\begin{solution}[0.7in]
		$\sigma_\beta = 1.5\mydeg = 0.0262\ \text{rad}$, so
		$y^2\sigma_\beta = 4(0.0262) = 0.105\ \m^2$, divided by $B$. Two
		neighboring stations fix depth four times worse than two stations
		$2\ \m$ apart.
		\end{solution}

	\part Bearings past $\pm 45\mydeg$ are poor: the beam is wide and the plot's
	edge can capture the peak marker. A box sits at the field's corner,
	$x = +1.0\ \m$. List the stations that see it within $\pm 45\mydeg$ of
	broadside when it is $1.0\ \m$ out and when it is $2.0\ \m$ out, and say
	why the tape runs $0.5\ \m$ past the field at each end.

	\begin{solution}[1.3in]
	A station sees the box within $\pm 45\mydeg$ when $|1.0 - x_i| \le y$.
	At $1.0\ \m$ out: stations $0$, $+0.5$, $+1.0$, $+1.5$. At $2.0\ \m$ out:
	$-1.0$ through $+1.5$, six stations. Without the station at $+1.5$, a near
	corner box would have only three useful stations and a short baseline.
	\end{solution}

	\part Write your station plan: the order in which you will visit the
	stations, and which two you take first. Give one sentence of reasoning.

	\begin{solution}[0.9in]
	Take the two end stations that hear the fox first, since they set the
	baseline and the depth error, then fill in the middle. If time runs short, the
	ends alone already give a good fix; two middle stations do not.
	\end{solution}

	\part Half the box spacing is $0.375\ \m$. Compare it with the depth error of
	a $2.0\ \m$ baseline for the deepest box, and say in one sentence whether
	naming the wrong box should be common.

	\begin{solution}[0.9in]
	$0.375/0.05 = 7.5$ standard deviations in depth, and more across range. With
	several good bearings, naming the wrong box is rare; a simulation of this
	range with a $2.5\mydeg$ budget, one bad bearing in seven, and a loosely aimed
	fox still names the right box about $93\%$ of the time. What goes wrong in
	practice is one bad bearing left in the solution.
	\end{solution}

\end{parts}

%%%%%%%%%%%%%%%%%%%%%%%%%%%%%%%%%%%%%%%%%%%%%%%%%%%
\newpage
\question \textbf{Your fox-hunt script.} Download \texttt{foxhunt\_starter.py}
from the L\labLesson\ Lab page. It reads \texttt{stations.csv} and
\texttt{boxes.csv} and prints the results; write its four functions in Python
with numpy. \texttt{solve} returns the least-squares fix and each station's
residual (Lesson 19, ``Many Bearings: Least Squares''); \texttt{drop\_worst}
finds the one station that disagrees with the rest; \texttt{rank\_boxes}
orders the boxes; and \texttt{amplitude\_check} scores how well the levels fit a
box. Test it on this practice set from an earlier class, recorded at one Rx Gain:

\begin{center}
\renewcommand{\arraystretch}{1.25}
\begin{tabular}{| c | c | c | c |}
\hline
\textbf{Station $x_i$ (m)} & \textbf{Angle $\theta_i$} & \textbf{Peak Gain} & \textbf{Use?} \\\hline\hline
$-1.5$ & $+52.9\mydeg$ & $-18.4$ dB & \kv{no} \\\hline
$-1.0$ & $+42.9\mydeg$ & $-17.3$ dB & \kv{yes} \\\hline
$-0.5$ & $+41.9\mydeg$ & $-14.3$ dB & \kv{yes} \\\hline
$0.0$  & $+18.6\mydeg$ & $-13.1$ dB & \kv{yes} \\\hline
$+0.5$ & $+2.9\mydeg$  & $-11.6$ dB & \kv{yes} \\\hline
$+1.0$ & $-16.5\mydeg$ & $-12.9$ dB & \kv{yes} \\\hline
$+1.5$ & $-30.4\mydeg$ & $-14.0$ dB & \kv{yes} \\\hline
\end{tabular}

\vspace{6pt}
Boxes: \textbf{1} $(-0.75,\ 1.20)$, \textbf{2} $(0.00,\ 1.00)$,
\textbf{3} $(0.55,\ 1.60)$, \textbf{4} $(-0.40,\ 1.90)$, \textbf{5} $(1.00,\ 0.90)$ m
\end{center}

\begin{parts}

	\part Mark in the table which stations you use, and say why you left any out.

	\begin{solution}[0.6in]
	Station $-1.5$ reads $52.9\mydeg$, past the $\pm 45\mydeg$ limit, so it is
	left out. The other six are inside it.
	\end{solution}

	\part Run \texttt{solve} on the stations you kept and record the fix and the
	residuals. Attach your code.

	\begin{center}
	\renewcommand{\arraystretch}{1.25}
	\begin{tabular}{| K{0.95in} | K{0.55in} | K{0.55in} | K{0.55in} | K{0.55in} | K{0.55in} | K{0.55in} |}
	\hline
	\textbf{Fix (m)} & $-1.0$ & $-0.5$ & $0.0$ & $+0.5$ & $+1.0$ & $+1.5$ \\\hline
	\kv{$(0.58,\ 1.49)$} & \kv{$0.14$} & \kv{$-0.20$} & \kv{$0.07$} & \kv{$0.00$} & \kv{$0.02$} & \kv{$-0.04$} \\\hline
	\end{tabular}
	\end{center}

		\begin{solution}[0.8in]
		Each station gives one equation in $(x, y)$,
		$x\cosp{\theta_i} - y\sinp{\theta_i} = x_i\cosp{\theta_i}$, scaled so the
		residual is each line's perpendicular distance from the fix. The signs
		depend on how the normal is chosen; the sizes are what matter.
		\end{solution}

	\part A good bearing here should leave a residual of a few centimeters, about
	$R\tanp{1.5\mydeg}$ at its range. Which station is the outlier? Run
	\texttt{drop\_worst}, and record the new fix and the largest residual left.

	\begin{center}
	\renewcommand{\arraystretch}{1.25}
	\begin{tabular}{| l | K{1.0in} | K{1.0in} | K{1.2in} |}
	\hline
	 & \textbf{Station dropped} & \textbf{Fix (m)} & \textbf{Largest residual left} \\\hline
	Outlier dropped & \kv{$-0.5$} & \kv{$(0.54,\ 1.63)$} & \kv{$0.04\ \m$} \\\hline
	\end{tabular}
	\end{center}

		\begin{solution}[1.2in]
		Station $-0.5$ has the largest residual, $0.20\ \m$, and its neighbor at
		$-1.0$ is also large because least squares spreads one bad line's error
		over the others. Dropping $-0.5$ leaves residuals of $0.04\ \m$ or less,
		which confirms it. That station read about $8\mydeg$ high, the size of
		error a reflection produces.
		\end{solution}

	\part Name the box from each fix and give its distance. Then compare the
	amplitude scores of the two nearest boxes. Does the level pattern agree
	with the bearings?

	\begin{solution}[1.2in]
	Both fixes are nearest box 3, $(0.55, 1.60)$: $0.11\ \m$ away with six
	stations and $0.03\ \m$ with the outlier dropped. The next nearest box is
	$0.73\ \m$ or more from either fix. An amplitude check that fits out the
	overall level and compares $1/R^2$ spreading with a patch roll-off near
	$\cos^{1.2}\theta$ fits box 3 within about $0.4\ \text{dB}$ RMS and every
	other box at $2.7\ \text{dB}$ or worse, so the levels agree. Scores depend on
	the model chosen; what counts is the ranking.
	\end{solution}

	\part The live box was box 3. Solve again with only the two end stations you
	kept ($-1.0$ and $+1.5$), then with only stations $0.0$ and $+0.5$. Record
	each fix, its distance from box 3, and the crossing angle of the pair.

	\begin{center}
	\renewcommand{\arraystretch}{1.25}
	\begin{tabular}{| l | K{1.0in} | K{1.0in} | K{1.0in} |}
	\hline
	\textbf{Pair} & \textbf{Fix (m)} & \textbf{To box 3} & \textbf{Crossing} \\\hline\hline
	$-1.0$ and $+1.5$ & \kv{$(0.53,\ 1.65)$} & \kv{$0.05\ \m$} & \kv{$73\mydeg$} \\\hline
	$0.0$ and $+0.5$ & \kv{$(0.59,\ 1.75)$} & \kv{$0.15\ \m$} & \kv{$16\mydeg$} \\\hline
	\end{tabular}
	\end{center}

		\begin{solution}[0.9in]
		The crossing angle is the difference of the two Angle readings, since every
		station shares one heading. The wide pair lands three times closer, and
		almost all of the narrow pair's miss is in depth, as Question 2(b) predicts.
		\end{solution}

\end{parts}

%%%%%%%%%%%%%%%%%%%%%%%%%%%%%%%%%%%%%%%%%%%%%%%%%%%
\newpage
\fullwidth{\large\textbf{Part 2: Measurements} \normalsize(your team, in class)}

Follow the procedure on the L\labLesson\ Lab page in read mode. Bring the kit up
as in the L18 Lab, with your own HB100 at boresight for \textbf{Find HB100} and
\textbf{Calibrate}, and set \textbf{Rx Gain (dB)} as the page says. Your own
HB100 stays on only for the reference shot. Then switch it off, disconnect its
power, and, once the instructor starts the hunt, press \textbf{Find HB100}
again so the kit tunes to the fox. Keep the calibration from bring-up. Once the
hunt starts, do not change Rx Gain.

\vspace{2pt}
\begin{center}
Own HB100 = \ansbox[1.2in]{$10.1$--$10.7\ \ghz$} \hspace{0.2in}
Fox, after re-tune = \ansbox[1.2in]{one value}\\[8pt]
Rx Gain for the hunt = \ansbox[1.2in]{one value, all stations}
\end{center}

%%%%%%%%%%%%%%%%%%%%%%%%%%%%%%%%%%%%%%%%%%%%%%%%%%%
\question \textbf{Reference shot.} Put the kit on station 0, squared to the tape.
Place your own HB100 $1\ \m$ downrange and $0.5\ \m$ toward the $+x$ end of the
tape, facing the kit, and record Angle. If it reads positive, $s = +1$; if
negative, $s = -1$. Multiply every Angle you record today by $s$.

\begin{center}
\renewcommand{\arraystretch}{1.25}
\begin{tabular}{| K{1.6in} | K{1.0in} |}
\hline
\textbf{Angle, source toward $+x$} & \textbf{$s$} \\\hline\hline
\kv{about $\pm 27\mydeg$} & \kv{$\pm 1$} \\\hline
\end{tabular}
\end{center}

\begin{solution}
$\tan^{-1}(0.5/1.0) = 26.6\mydeg$. The L18 Lab put positive Angle on the E8
end of the array, so $s$ depends only on which way round the board sits on its
stand. A reading more than $5\mydeg$ from $26.6\mydeg$ in size means the kit is
not square to the tape or the source is not where you think.
\end{solution}

%%%%%%%%%%%%%%%%%%%%%%%%%%%%%%%%%%%%%%%%%%%%%%%%%%%
\question \textbf{The hunt.} At each station, square the kit to the tape with the
middle of the array on the mark, and read Angle on ten consecutive sweeps
from the \textbf{Tracking} tab. Record $s$ times the mean, the spread (largest
minus smallest), and Peak Gain. Write ``none'' for a station with no clear
tone. Ends first.

\begin{center}
\small
\renewcommand{\arraystretch}{1.25}
\begin{tabular}{| c | K{1.1in} | K{0.8in} | K{1.0in} | K{0.7in} |}
\hline
\textbf{Station $x_i$ (m)} & \textbf{$s\,\bar\theta_i$} & \textbf{Spread} & \textbf{Peak Gain} & \textbf{Use?} \\\hline\hline
$-1.5$ & \kv{} & \kv{$\lesssim 2\mydeg$} & \kv{} & \kv{} \\\hline
$-1.0$ & \kv{} & \kv{$\lesssim 2\mydeg$} & \kv{} & \kv{} \\\hline
$-0.5$ & \kv{} & \kv{$\lesssim 2\mydeg$} & \kv{} & \kv{} \\\hline
$0.0$  & \kv{} & \kv{$\lesssim 2\mydeg$} & \kv{} & \kv{} \\\hline
$+0.5$ & \kv{} & \kv{$\lesssim 2\mydeg$} & \kv{} & \kv{} \\\hline
$+1.0$ & \kv{} & \kv{$\lesssim 2\mydeg$} & \kv{} & \kv{} \\\hline
$+1.5$ & \kv{} & \kv{$\lesssim 2\mydeg$} & \kv{} & \kv{} \\\hline
\end{tabular}
\end{center}

\begin{solution}
Every row depends on the room. Check that each angle used is within
$\pm 45\mydeg$, that the spread is a degree or two (near $5\mydeg$ suggests the
receiver is near full scale), and that the angles fall steadily from one end of
the line to the other, as they must for one source downrange.
\end{solution}

%%%%%%%%%%%%%%%%%%%%%%%%%%%%%%%%%%%%%%%%%%%%%%%%%%%
\question \textbf{The call.} Run your Part 1 script on your stations in class,
before the reveal. Record the fix, the largest residual, any station dropped, the
box you name, and, after the instructor opens the boxes, the box that was live.

\begin{center}
\renewcommand{\arraystretch}{1.25}
\begin{tabular}{| K{1.1in} | K{1.0in} | K{0.9in} | K{0.9in} | K{0.9in} |}
\hline
\textbf{Fix (m)} & \textbf{Largest residual} & \textbf{Dropped} & \textbf{Box named} & \textbf{Live box} \\\hline\hline
\kv{} & \kv{$\lesssim 0.15\ \m$} & \kv{} & \kv{} & \kv{} \\\hline
\end{tabular}
\end{center}

%%%%%%%%%%%%%%%%%%%%%%%%%%%%%%%%%%%%%%%%%%%%%%%%%%%
\newpage
\fullwidth{\large\textbf{Part 3: Writeup} \normalsize(your team, after the lab)}

%%%%%%%%%%%%%%%%%%%%%%%%%%%%%%%%%%%%%%%%%%%%%%%%%%%
\question \textbf{Reconcile the fix.} The live box's measured position is posted
after the period. Answer each part in three or four sentences, using your numbers
from Part 2 and your predictions from Part 1.

\begin{parts}

	\part Give the distance from your fix to the live box, split into its
	across-range and depth parts. Compare each with your Question 2 predictions.
	\begin{solution}[1.5in]
	With a long baseline the miss should be a few centimeters across range and
	somewhat more in depth, together $0.15$ to $0.2\ \m$ or less. A depth miss
	much larger than across range means the baseline was short or the end
	stations were dropped; a large across-range miss points at a wrong $s$, a kit
	not squared to the tape, or a kit below the boxes.
	\end{solution}

	\part Use the live box's position to find the true angle from each station,
	$\tan^{-1}\left((x_\text{box} - x_i)/y_\text{box}\right)$, and the error of
	each measured angle. Which term of your error budget best explains the largest
	one?
	\begin{solution}[1.8in]
	Errors of $1$ to $2\mydeg$ are within the budget. Errors that share one
	sign and size at every station point at the heading common to all of them,
	the kit not square to the tape. One large error at a single station points at
	a reflection, especially if that station's spread was large or its level out
	of line with its neighbors.
	\end{solution}

	\part If your script dropped a station, say whether dropping it moved the fix
	toward the live box, and what in the room you think caused it. If it dropped
	none, say what that tells you about your bearings.
	\begin{solution}[1.5in]
	A real outlier is confirmed when dropping it shrinks the other residuals and
	moves the fix toward the box. Whiteboards, metal cabinets, and people standing
	near a line of sight are the usual causes. If every residual stayed near
	$R_i\tanp{\sigma_\beta}$, the bearings agree within the budget.
	\end{solution}

\end{parts}

%%%%%%%%%%%%%%%%%%%%%%%%%%%%%%%%%%%%%%%%%%%%%%%%%%%
\newpage
\question \textbf{Baseline and amplitude.}

\begin{parts}

	\part Solve your own data again with only your two end stations, then with
	only two neighboring stations near the middle. Record both fixes and their
	distances to the live box. Does the result match Question 2(b)?
	\begin{solution}[1.8in]
	The end pair should land close, with most of any miss in depth; the
	neighboring pair should miss mostly in depth, by several times more, roughly
	in the ratio of the two baselines. A neighboring pair that happens to land
	close has errors that happened to cancel; a two-station fix has no residuals,
	so it cannot tell you which case you are in.
	\end{solution}

	\part In two or three sentences, explain why Peak Gain can help choose
	between two boxes but cannot give the fox's range on its own.
	\begin{solution}[1.3in]
	The level depends on range, but also on which way the fox faces, on
	reflections, and on the receiver's gain, none of which you know well enough
	to turn one level into a distance. Comparing the pattern of levels along the
	line, with the overall level fitted out, removes the unknown gain and much of
	the fox's aim, so the shape can still separate two candidates that the
	bearings place about equally close.
	\end{solution}

\end{parts}

\end{questions}

\end{document}
